\chapter{Pipeline d'extraction du graphe métier}

\section{Motivation et modèle}

Les fiches expriment des connaissances fonctionnelles qui ne figurent pas toujours comme relations UML. La couche métier introduit des nœuds \code{ConceptMetier} séparés des classes. Les principaux liens sont :

\begin{center}
\resizebox{0.92\linewidth}{!}{%
\begin{tikzpicture}[node distance=2cm and 1.8cm]
  \node[stockage] (cm) {ConceptMetier\\Poste d'aiguillage};
  \node[source,right=of cm] (cl) {Classe UML\\Poste};
  \node[source,below=of cm] (doc) {Document\\Post.txt};
  \node[source,right=of doc] (chunk) {Chunk\\passage 1};
  \draw[fleche] (cm)--node[above,font=\scriptsize]{REPRESENTE\_PAR}(cl);
  \draw[fleche] (cm)--node[left,font=\scriptsize]{DOCUMENTE\_PAR}(doc);
  \draw[fleche] (cm)--node[above,sloped,font=\scriptsize]{ETAYE\_PAR}(chunk);
\end{tikzpicture}
}
\end{center}

\section{Sélection déterministe des classes candidates}

Le LLM ne reçoit pas les 998 classes. Pour chaque passage, la fonction \code{classes\_candidates} charge le catalogue JSON et normalise le texte. Une classe reçoit un score de 1 si son nom est présent exactement. Sinon, le score correspond au recouvrement entre les mots du passage et ceux du nom de la classe :

\begin{equation}
score(c,p)=\frac{|mots(c)\cap mots(p)|}{|mots(c)|}.
\end{equation}

Les classes dont le score atteint 0,5 sont triées et les 30 premières sont transmises à Mistral. Pour chacune, le prompt fournit le nom, jusqu'à 300 caractères de description et 15 attributs. La sélection actuelle est lexicale ; elle pourrait être améliorée par une segmentation systématique du CamelCase et par des synonymes contrôlés.

\section{Extraction par le LLM}

Le passage, ses lignes et les classes candidates sont insérés dans un prompt d'extraction. Mistral doit retourner exclusivement un JSON contenant :

\begin{itemize}
  \item les concepts métier et leurs alias ;
  \item une catégorie parmi équipement, objet, type, acteur, document, procédure, localisation, date, valeur, caractéristique ou système ;
  \item les correspondances exactes vers les classes candidates ;
  \item des triplets source--prédicat--cible ;
  \item une citation exacte pour chaque relation.
\end{itemize}

Il ne s'agit pas d'un modèle NER spécialisé : les entités sont identifiées par les capacités générales du LLM, guidées par le schéma de sortie et des exemples positifs et négatifs.

\begin{lstlisting}[language={},caption={Sortie attendue de l'extracteur}]
{
  "concepts": [
    {"id_local": "poste_aiguillage",
     "nom": "Poste d'aiguillage",
     "categorie": "EQUIPEMENT",
     "classes_arianes": ["Poste"]}
  ],
  "relations": [
    {"source": "signal",
     "type": "EST_COMMANDE_PAR",
     "cible": "poste_aiguillage",
     "citation_exacte": "..."}
  ]
}
\end{lstlisting}

\section{Contrôles déterministes}

La sortie n'est jamais écrite directement dans Neo4j. Le code vérifie successivement :

\begin{enumerate}
  \item la validité du JSON et des identifiants normalisés ;
  \item l'existence exacte des classes Arianes proposées ;
  \item l'existence des concepts source et cible ;
  \item l'absence d'auto-relation ;
  \item l'appartenance du prédicat à une liste contrôlée ;
  \item la présence textuelle de la citation dans le passage ;
  \item la compatibilité de certaines catégories avec le prédicat ;
  \item la présence d'indices lexicaux justifiant le prédicat ;
  \item les négations et modalités telles que « peut être » ou « muni ou non ».
\end{enumerate}

Une absence d'indice lexical produit un avertissement plutôt qu'un rejet automatique. La validation V3 classe ensuite la relation comme \code{VALIDE}, \code{PROBABLE} ou \code{REJETE} selon ses preuves et son score.

\section{Direction et corrections bornées}

La direction d'un triplet est stricte : « A commande B » donne \code{A COMMANDE B}, tandis que « B est commandé par A » donne \code{B EST\_COMMANDE\_PAR A}. Pour \code{EST\_UN}, le type spécialisé est la source. Les inversions ne sont corrigées que pour des cas relus et accompagnés de la citation attendue. Chaque correction conserve l'ancien triplet et son motif.

\section{Consolidation}

Les concepts extraits de plusieurs passages sont fusionnés par nom normalisé. Les alias, classes et documents sont réunis. Une incompatibilité de catégories conduit à \code{INCONNU}. Les extrémités des relations sont ensuite réécrites vers les identifiants canoniques et les auto-relations sont contrôlées une seconde fois. Deux preuves différentes d'un même triplet restent distinctes.

\section{Cycle des statuts}

\begin{center}
\resizebox{0.92\linewidth}{!}{%
\begin{tikzpicture}[node distance=2cm]
  \node[source] (llm) {Proposition LLM};
  \node[traitement,right=of llm] (ctrl) {Contrôles};
  \node[traitement,right=of ctrl] (review) {PROBABLE};
  \node[stockage,above right=1cm and 1.6cm of review] (valid) {VALIDE documentaire};
  \node[bloc,below right=1cm and 1.6cm of review,fill=sncfrouge!8] (reject) {REJETE};
  \draw[fleche] (llm)--(ctrl); \draw[fleche] (ctrl)--(review);
  \draw[fleche] (review)--node[above,sloped,font=\scriptsize]{relecture}(valid);
  \draw[fleche] (ctrl)--(reject); \draw[fleche] (review)--(reject);
\end{tikzpicture}
}
\end{center}

Les objets \code{VALIDE} et \code{PROBABLE} sans motif bloquant peuvent être importés. Leur statut reste attaché à la relation et contrôle la prudence de la réponse.

\section{État réel du corpus}

Les 29 fiches ont été traitées. Le recontrôle V3 a observé 345 concepts, retenu 117 relations \code{VALIDE} et 224 relations \code{PROBABLE}, et isolé 24 relations \code{REJETE}. Trente-six relations retenues ont bénéficié d'une correction bornée.

\attention{Le pilote documentaire global est désormais importé. La base observée contient 255 concepts métier et 347 relations actives : 123 \code{VALIDE} et 224 \code{PROBABLE}. Le statut \code{VALIDE} reste documentaire et non expert.}

\capture{neo4j-pilote-poste.png}{Exemple de sous-graphe métier autour du poste d'aiguillage}{fig:pilote-poste}
